\section{Beyond a Scalar Effective Sample Size}
\label{sec:ess}

Theorem~\ref{thm:scaling} suggests a mode-wise replacement for the usual
scalar effective sample size.  After $B$ refreshes, mode $j$ has expected
innovation count
\begin{equation}
Bq_j
=\frac{B}{Z}\frac{\lambda_j}{\tau_j}.
\label{eq:mode-wise-information}
\end{equation}
Thus $\lambda_j$ measures how much test risk lies in a direction, whereas
$\lambda_j/\tau_j$ measures how rapidly independent information about that
direction arrives.  A mode becomes statistically visible near $Bq_j\asymp1$.
The full vector $(Bq_j)_j$, rather than one global count, determines the
learnable spectral resolution.

\paragraph{A matched-autocorrelation counterexample.}
The distinction can be made formal using the reversible Markov realization of
the supplement.  Its lag-$\ell$ covariance is
\begin{equation}
C(\ell)=\sum_{j\ge1}\lambda_j
(1-\tau_j^{-1})^\ell e_je_j^\top.
\label{eq:mode-covariance-main}
\end{equation}
A natural global summary is the trace integrated autocorrelation time
\begin{equation}
\tau_{\rm tr}
:=1+2\sum_{\ell\ge1}
\frac{\operatorname{tr}C(\ell)}{\operatorname{tr}C(0)}
=2\sum_j\lambda_j\tau_j-1.
\label{eq:trace-iat}
\end{equation}
It is the variance-inflation factor obtained by averaging autocorrelation over
the marginal spectral mass, analogous to the integrated autocorrelation time
used in time-series and Markov-chain Monte Carlo diagnostics
\cite{geyer1992practical}.  The next proposition shows that even this
well-matched scalar does not determine the learning exponent.

\begin{proposition}[Same global autocorrelation, different scaling]
\label{prop:matched-iat}
Let $\lambda_j=j^{-a}/\zeta(a)$ and $s_j=j^{-b}/\zeta(b)$, where $a>1$ and
$1<b<a+r+1$.  Fix $0<r<a-1$, define
$m_r=\sum_j\lambda_jj^r<\infty$, and choose $T\ge m_r$.  Compare
\begin{equation}
\tau_j^{\rm unif}=T,
\qquad
\tau_j^{\rm align}=\frac{T}{m_r}j^r.
\label{eq:matched-profiles}
\end{equation}
The two stationary streams have the same marginal covariance and the same
trace autocorrelation $\tau_{\rm tr}=2T-1$.  Nevertheless, their innovation
laws are
\begin{equation}
q_j^{\rm unif}=\frac{j^{-a}}{\zeta(a)},
\qquad
q_j^{\rm align}=\frac{j^{-(a+r)}}{\zeta(a+r)},
\label{eq:matched-innovation-laws}
\end{equation}
and, when model truncation is nonbinding, their noiseless risks scale as
$B^{-(b-1)/a}$ and $B^{-(b-1)/(a+r)}$, respectively.  The same two exponents
hold for stationary raw horizons $N$.
\end{proposition}

\begin{proof}
Equation~\eqref{eq:trace-iat} gives $2T-1$ for both profiles because
$\sum_j\lambda_j(T/m_r)j^r=T$.  Normalizing
$\lambda_j/\tau_j$ yields Equation~\eqref{eq:matched-innovation-laws}, and
Theorem~\ref{thm:scaling} gives the two block exponents.  For raw horizons,
Theorem~\ref{thm:raw} applies.  Since $r<a-1$, the equilibrium residual-life
exponent $(a-1)/r$ exceeds one, while $(b-1)/(a+r)<1$, so the inspection term
is asymptotically faster and does not change the aligned exponent.
\end{proof}

The proposition is deliberately specific: it does not rule out every possible
scalar encoding of an entire process.  It does rule out the common practice of
using a single global autocorrelation correction while ignoring how temporal
persistence aligns with predictive spectral directions.  Two streams can
receive exactly the same scalar correction $N/(2T-1)$ and still have different
power-law slopes.  Figure~\ref{fig:mechanism-planning} verifies this
construction directly.

\begin{figure*}[t]
\centering
\includegraphics[width=\textwidth]{mechanism_planning.pdf}
\caption{Why a scalar effective sample size is insufficient and how the
replacement affects planning. (a) Uniform and spectrally aligned streams have
the same marginal spectrum and are calibrated to the same trace integrated
autocorrelation time, but the aligned innovation spectrum decays faster.
(b) Stationary raw-horizon simulations recover different slopes despite the
matched global autocorrelation. (c) Under dense-training cost $C=MB$,
increasing persistence shifts the optimum toward more innovations and reduces
the compute exponent.}
\label{fig:mechanism-planning}
\end{figure*}

\paragraph{Data and compute planning.}
The exponent has a direct sample-complexity interpretation.  In the noiseless
data-limited regime, reaching risk $\eps$ requires
\begin{equation}
B_\eps\asymp \eps^{-(a+r)/(b-1)}.
\label{eq:sample-complexity}
\end{equation}
Relative to uniform persistence, aligned persistence therefore multiplies the
required number of innovations by
$\eps^{-r/(b-1)}$.  The penalty grows as the desired error becomes smaller,
which is why a constant-factor correction cannot repair the slope.

For a stylized dense learner whose training cost is proportional to
$C=MB$, the model and innovation terms can also be optimized jointly.

\begin{corollary}[Compute-optimal frontier]
\label{cor:compute-optimal}
Let $p=a+r$ and suppose the noiseless power-law assumptions of
Theorem~\ref{thm:scaling} hold.  Under $MB\le C$,
\begin{align}
M_\star&\asymp C^{1/(p+1)},
&B_\star&\asymp C^{p/(p+1)}
\label{eq:compute-allocation}\\
\inf_{MB\le C}\E\risk
&\asymp C^{-(b-1)/(p+1)}
\label{eq:compute-rate}.
\end{align}
\end{corollary}

\begin{proof}
Choosing the allocations in Equation~\eqref{eq:compute-allocation} balances
$M^{-(b-1)}$ and $B^{-(b-1)/p}$ and gives the upper bound.  Conversely, if
$MB\le C$, then at least one of
$M\lesssim C^{1/(p+1)}$ or $B\lesssim C^{p/(p+1)}$ must hold.  The
corresponding term in Theorem~\ref{thm:scaling} gives the matching lower bound.
\end{proof}

Persistence therefore changes not only the achievable risk but the allocation
rule: as $r$ grows, the optimal model exponent $1/(a+r+1)$ decreases and the
innovation exponent $(a+r)/(a+r+1)$ increases.  Under this cost model, more of
the budget should be spent obtaining distinct temporal or environmental
coverage rather than enlarging a model that cannot access the missing modes.

\begin{table}[t]
\centering
\footnotesize
\caption{Dominant asymptotic regimes, with $p=a+r$.  The full expressions and
conditions are given in Theorems~\ref{thm:scaling}--\ref{thm:noisy-raw}.}
\label{tab:regimes}
\setlength{\tabcolsep}{1.5pt}%
\begin{tabular}{@{}lp{0.33\columnwidth}p{0.29\columnwidth}@{}}
\toprule
Regime & Characteristic choice or condition & Dominant risk \\
\midrule
Model limited & $M\ll B^{1/p}$ & $M^{-(b-1)}$ \\
Innovation limited & $M\gtrsim B^{1/p}$, $b\le p$ & $B^{-(b-1)/p}$ \\
Noise limited & $b>p$, $M\asymp(B/\sigma^2)^{1/b}$ & $(\sigma^2/B)^{(b-1)/b}$ \\
Compute optimal & $C=MB$ & $C^{-(b-1)/(p+1)}$ \\
Stationary heavy dwell & $(a-1)/r<(b-1)/p$ & $N^{-(a-1)/r}$ \\
\bottomrule
\end{tabular}
\end{table}
